% GENERATED FILE. Edit canonical lessons, then run npm run generate:books.
% canonical-source-hash: fnv1a64:6c0421b02855fda2
% canonical-lessons: GU-C146-shun-polar, GU-C146-ke-nahi, GU-C146-mane, GU-C146-maaro, GU-C146-tamaro, GU-C146-ask-practice

\chapter{Asking yes or no}
\label{ch:gu-asking-yes-or-no}
\hlchaptermodality{146}

\begin{chapteropening}
\textbf{By the end of this chapter:} \emph{I can ask a question that is answered yes or no, with śũ, ke nahī̃ or kharũ, answer it, and say whose things are whose and what someone likes or has.}

Carry a short exchange of yes-or-no questions about likes, family and plans, and answer one with a negative.
\end{chapteropening}

\section[\texorpdfstring{\gu{શું} …? (śũ …?)}{શું …? (śũ …?)}]{\gu{શું} …? — the yes-or-no question}

\label{lesson:GU-C146-shun-polar}



\begin{quote}
Say \textbf{\gu{શું}}, \emph{what}, as in \emph{your name}: \textbf{\gu{તમારું} \gu{નામ} \gu{શું} \gu{છે}}. Then the two oldest answers in the book, \textbf{\gu{હા}} and \textbf{\gu{ના}}. And one sentence from the last chapter: \emph{I don't speak Gujarati}.
\end{quote}

\begin{grammarlens}[title={Grammar Lens: \gu{શું} at the front}]
\begin{quote}
\textbf{\gu{શું} \gu{તમે} \gu{ગુજરાતી} \gu{બોલો} \gu{છો}?} — \emph{śũ tame gujarātī bolo cho?} — Do you speak Gujarati?
\end{quote}

Inside a sentence \textbf{\gu{શું}} means \emph{what}. At the very front it means nothing at all: it marks the sentence as a yes-or-no question, and not one other word moves. English rebuilds the sentence around \emph{do}; Gujarati sets a flag at the front.

\begin{itemize}
\raggedright
  \item \textbf{\gu{તમારું} \gu{નામ} \gu{શું} \gu{છે}} — \emph{what} is your name — inside;
  \item \textbf{\gu{શું} \gu{તમે} \gu{ગુજરાતી} \gu{બોલો} \gu{છો}} — \emph{do} you speak Gujarati — in front.
\end{itemize}
\end{grammarlens}

\subsection*{You'll want to know: the voice can carry it}
In speech the flag often drops, and the voice does its work: \textbf{\gu{તમે} \gu{ગુજરાતી} \gu{બોલો} \gu{છો}?} with the pitch rising on the last word asks the same question. On the page the question mark carries the rise.

\begin{grammarlens}[title={Grammar Lens: answering}]
Gujarati answers with \textbf{\gu{હા}} or \textbf{\gu{ના}}, and then, usually, with the verb again:

\begin{quote}
\textbf{\gu{હા}, \gu{હું} \gu{ગુજરાતી} \gu{બોલું} \gu{છું}.} — \emph{hā, hũ gujarātī bolũ chũ.} \textbf{\gu{ના}, \gu{હું} \gu{ગુજરાતી} \gu{નથી} \gu{બોલતો}.} — \emph{nā, hũ gujarātī nathī bolto.}
\end{quote}
\end{grammarlens}

\subsection*{Your turn}
\begin{itemize}
\raggedright
  \item \emph{Say it:} \textbf{\gu{શું} \gu{તમે} \gu{અહીં} \gu{રહો} \gu{છો}?} — \emph{śũ tame ahī̃ raho cho?}, do you live here?
  \item \emph{Answer:} with \textbf{\gu{હા}} and the verb again
  \item \emph{Say it:} the same question with no \textbf{\gu{શું}}, the voice rising at the end
\end{itemize}

\begin{tcolorbox}[breakable,colback=teal!4,colframe=teal!35!black,title={Before you move on}]
What does \textbf{\gu{શું}} mean at the very front? (\textbf{Nothing}; it marks a yes-or-no question.) What else moves? (\textbf{Nothing}.)
\end{tcolorbox}
\section[\texorpdfstring{… \gu{કે} \gu{નહીં}? · …, \gu{ખરું}? (… ke nahī̃? · … …)}{… કે નહીં? · …, ખરું? (… ke nahī̃? · … …)}]{\gu{કે} \gu{નહીં}? and \gu{ખરું}? — two tags at the end}

\label{lesson:GU-C146-ke-nahi}



\begin{quote}
Ask \emph{do you speak Gujarati?} twice: once with \textbf{\gu{શું}}, once with your voice alone. Then say \textbf{\gu{નહીં}}, the \emph{not} for one word, and \textbf{\gu{સાચું}}, \emph{true}.
\end{quote}

\begin{grammarlens}[title={Grammar Lens: … \gu{કે} \gu{નહીં}?}]
\begin{quote}
\textbf{\gu{તમે} \gu{આવો} \gu{છો} \gu{કે} \gu{નહીં}?} — \emph{tame āvo cho ke nahī̃?} — Are you coming or not?
\end{quote}

The \textbf{\gu{કે}} you know as \emph{that} also means \emph{or} between two choices. The tag offers both answers and asks the listener to pick one. It is firmer than \textbf{\gu{શું}}: a parent at the door, a friend whose bus is leaving.
\end{grammarlens}

\subsection*{You'll want to know: …, \gu{ખરું}?}
\textbf{\gu{ખરું}} (\emph{kharũ}) means \emph{true, real}, a close neighbour of \textbf{\gu{સાચું}}. As a tag it asks for agreement:

\begin{quote}
\textbf{\gu{તમે} \gu{અહીં} \gu{રહો} \gu{છો}, \gu{ખરું}?} — \emph{tame ahī̃ raho cho, kharũ?} — You live here, right?
\end{quote}

The speaker already expects \textbf{\gu{હા}} and checks it. So the three sound different: \textbf{\gu{શું}} in front opens the question, \textbf{\gu{કે} \gu{નહીં}} at the end asks for a pick, and \textbf{\gu{ખરું}} at the end asks for a yes.

\subsection*{Your turn}
\begin{itemize}
\raggedright
  \item \emph{Say it:} \textbf{\gu{તમે} \gu{ચા} \gu{પીઓ} \gu{છો} \gu{કે} \gu{નહીં}?} — \emph{tame cā pīo cho ke nahī̃?}
  \item \emph{Say it:} \textbf{\gu{તે} \gu{ગુજરાતી} \gu{બોલે} \gu{છે}, \gu{ખરું}?} — \emph{te gujarātī bole che, kharũ?}
  \item \emph{Say it:} \textbf{\gu{આપણે} \gu{ચા} \gu{પીએ} \gu{છીએ}, \gu{ખરું}?} — \emph{āpṇe cā pīe chīe, kharũ?}, we drink tea, right?
  \item \emph{Answer:} each one, with \textbf{\gu{હા}} or \textbf{\gu{ના}} and the verb again
\end{itemize}

\begin{tcolorbox}[breakable,colback=teal!4,colframe=teal!35!black,title={Before you move on}]
Which tag already expects \emph{yes}? (\textbf{\gu{ખરું}?}) What does \textbf{\gu{કે}} mean in \textbf{\gu{કે} \gu{નહીં}?} (\emph{Or}.)
\end{tcolorbox}
\section[\texorpdfstring{\gu{મને} · \gu{તને} · \gu{તેને} (mane · tane · tene)}{મને · તને · તેને (mane · tane · tene)}]{\gu{મને}, \gu{તને}, \gu{તેને} — the pronouns things happen to}

\label{lesson:GU-C146-mane}



\begin{quote}
Ask \emph{are you coming or not?}, then \emph{you live here, right?} And from the verb chapters long ago: \textbf{\gu{મને} \gu{ગુજરાતી} \gu{ગમે} \gu{છે}}, \emph{I like Gujarati}.
\end{quote}

\begin{grammarlens}[title={Grammar Lens: the person something happens to}]
In \textbf{\gu{મને} \gu{ગુજરાતી} \gu{ગમે} \gu{છે}} the thing liked is the subject, and \emph{I} stand to one side as \textbf{\gu{મને}}, \emph{to me}. Gujarati speaks of a fever the same way: it arrives \emph{to} someone.

\begin{quote}
\textbf{\gu{મને} \gu{તાવ} \gu{છે}.} — \emph{mane tāv che.} — I have a fever.
\end{quote}

The \emph{to} forms for one person:

\noindent
\begin{tabularx}{\linewidth}{@{}>{\raggedright\arraybackslash}X>{\raggedright\arraybackslash}X>{\raggedright\arraybackslash}X@{}}
\toprule
\textbf{who} & \textbf{pronoun} & \textbf{\emph{to} form} \\
\midrule
I & \gu{હું} & \gu{મને} \\
you (familiar) & \gu{તું} & \gu{તને} \\
he, she & \gu{તે} & \gu{તેને} \\
\bottomrule
\end{tabularx}
\end{grammarlens}

\begin{grammarlens}[title={Grammar Lens: more than one}]
\textbf{\gu{તમને}} has been in the book since the second day, inside \textbf{\gu{તમને} \gu{મળીને} \gu{આનંદ} \gu{થયો}}: \emph{to you, meeting, joy happened}. Its partners:

\noindent
\begin{tabularx}{\linewidth}{@{}>{\raggedright\arraybackslash}X>{\raggedright\arraybackslash}X>{\raggedright\arraybackslash}X@{}}
\toprule
\textbf{who} & \textbf{pronoun} & \textbf{\emph{to} form} \\
\midrule
we & \gu{અમે} & \gu{અમને} \\
you (respectful) & \gu{તમે} & \gu{તમને} \\
they & \gu{તેઓ} & \gu{તેમને} \\
\bottomrule
\end{tabularx}

\textbf{\gu{તેઓ}} turns into \textbf{\gu{તેમ}-} before the ending: \textbf{\gu{તેમને}}, \emph{to them}.
\end{grammarlens}

\subsection*{Your turn}
Say these aloud:

\begin{itemize}
\raggedright
  \item \textbf{\gu{મને} \gu{ચા} \gu{ગમે} \gu{છે}} — \emph{mane cā game che}
  \item \textbf{\gu{શું} \gu{તમને} \gu{તાવ} \gu{છે}?} — \emph{śũ tamne tāv che?}, do you have a fever?
  \item \textbf{\gu{તેમને} \gu{ગુજરાતી} \gu{ગમે} \gu{છે}} — \emph{temne gujarātī game che}, they like Gujarati
\end{itemize}

\begin{tcolorbox}[breakable,colback=teal!4,colframe=teal!35!black,title={Before you move on}]
Who is the subject of \textbf{\gu{મને} \gu{ચા} \gu{ગમે} \gu{છે}}? (The tea; \emph{I} am \textbf{\gu{મને}}, \emph{to me}.) What does \textbf{\gu{તેઓ}} become before \emph{to}? (\textbf{\gu{તેમને}}.)
\end{tcolorbox}
\section[\texorpdfstring{\gu{મારો} · \gu{મારી} … (māro · mārī …)}{મારો · મારી … (māro · mārī …)}]{\gu{મારો} \gu{ભાઈ}, \gu{મારી} \gu{બહેન} — \emph{my} agrees with what is owned}

\label{lesson:GU-C146-maaro}



\begin{quote}
Say \emph{I have a fever} and \emph{they like Gujarati}. Then the first possessive you ever learned: \textbf{\gu{મારું} \gu{નામ}}, \emph{my name}.
\end{quote}

\begin{grammarlens}[title={Grammar Lens: the ending follows the thing}]
\textbf{\gu{મારું}} ends in \textbf{-\gu{ું}} because \textbf{\gu{નામ}} is a neuter noun. Change the noun and \emph{my} changes its last sign to match it:

\begin{itemize}
\raggedright
  \item \textbf{\gu{મારો} \gu{ભાઈ}} — \emph{māro bhāī} — my brother: masculine, \textbf{-\gu{ો}};
  \item \textbf{\gu{મારી} \gu{બહેન}} — \emph{mārī bahen} — my sister: feminine, \textbf{-\gu{ી}};
  \item \textbf{\gu{મારું} \gu{ઘર}} — \emph{mārũ ghar} — my home: neuter, \textbf{-\gu{ું}};
  \item \textbf{\gu{મારા} \gu{પિતા}} — \emph{mārā pitā} — my father: \textbf{-\gu{ા}}, for more than one, or for one person spoken of with respect.
\end{itemize}

A woman and a man say \textbf{\gu{મારો} \gu{ભાઈ}} alike. The ending follows the brother, not whoever owns him.
\end{grammarlens}

\subsection*{Your turn}
Say these aloud:

\begin{itemize}
\raggedright
  \item \textbf{\gu{મારા} \gu{પિતા} \gu{અહીં} \gu{છે}} — \emph{mārā pitā ahī̃ che}
  \item \textbf{\gu{મારી} \gu{ચા}} — \emph{mārī cā}, my tea, and why it is \textbf{-\gu{ી}} (\textbf{\gu{ચા}} is feminine)
\end{itemize}

\begin{tcolorbox}[breakable,colback=teal!4,colframe=teal!35!black,title={Before you move on}]
Whose gender decides the ending of \emph{my}? (The thing owned.) Why \textbf{\gu{મારા} \gu{પિતા}}? (Respect takes \textbf{-\gu{ા}}.)
\end{tcolorbox}
\section[\texorpdfstring{\gu{તારું} · \gu{તેનું} · \gu{અમારું} … (tārũ · tenũ …)}{તારું · તેનું · અમારું … (tārũ · tenũ …)}]{\gu{તમારો} \gu{ફોન}, \gu{તેની} \gu{બહેન} — every possessive bends the same way}

\label{lesson:GU-C146-tamaro}



\begin{quote}
Say \emph{my brother}, \emph{my sister} and \emph{my home}. Then the first question you ever asked: \textbf{\gu{તમારું} \gu{નામ} \gu{શું} \gu{છે}?}
\end{quote}

\begin{grammarlens}[title={Grammar Lens: one owner}]
\noindent
\begin{tabularx}{\linewidth}{@{}>{\raggedright\arraybackslash}X>{\raggedright\arraybackslash}X@{}}
\toprule
\textbf{owner} & \textbf{\emph{my}, \emph{your}, \emph{his}: the -\gu{ું} form} \\
\midrule
\gu{હું} & \gu{મારું} \\
\gu{તું} & \gu{તારું} \\
\gu{તે} & \gu{તેનું} \\
\bottomrule
\end{tabularx}

\begin{quote}
\textbf{\gu{તેનું} \gu{નામ} \gu{મીરા} \gu{છે}.} — \emph{tenũ nām mīrā che.} — Her name is Mira.
\end{quote}
\end{grammarlens}

\begin{grammarlens}[title={Grammar Lens: more than one owner}]
\noindent
\begin{tabularx}{\linewidth}{@{}>{\raggedright\arraybackslash}X>{\raggedright\arraybackslash}X@{}}
\toprule
\textbf{owner} & \textbf{the -\gu{ું} form} \\
\midrule
\gu{અમે} & \gu{અમારું} \\
\gu{આપણે} & \gu{આપણું} \\
\gu{તમે} & \gu{તમારું} \\
\gu{તેઓ} & \gu{તેમનું} \\
\bottomrule
\end{tabularx}

\textbf{\gu{તેમનું}} comes from the same \textbf{\gu{તેમ}-} as \textbf{\gu{તેમને}}. Every one of them bends exactly as \textbf{\gu{મારું}} does, to \textbf{-\gu{ો}}, \textbf{-\gu{ી}}, \textbf{-\gu{ું}} or \textbf{-\gu{ા}}:

\begin{itemize}
\raggedright
  \item \textbf{\gu{તમારો} \gu{ફોન}} — \emph{tamāro phon} — your phone;
  \item \textbf{\gu{તેની} \gu{બહેન}} — \emph{tenī bahen} — his sister, or her sister;
  \item \textbf{\gu{આપણું} \gu{ગામ}} — \emph{āpṇũ gām} — our village, yours and mine;
  \item \textbf{\gu{તમારા} \gu{પિતા}}, \textbf{\gu{તેના} \gu{પિતા}} — \emph{tamārā pitā}, \emph{tenā pitā} — your father, his or her father: respect takes \textbf{-\gu{ા}}.
\end{itemize}
\end{grammarlens}

\subsection*{Your turn}
Say these aloud:

\begin{itemize}
\raggedright
  \item \textbf{\gu{તારું} \gu{નામ} \gu{શું} \gu{છે}?} — \emph{tārũ nām śũ che?}, to a child
  \item \emph{amārũ ghar}, then \emph{āpṇũ ghar}, and who is left out of the first
  \item \textbf{\gu{તેમનો} \gu{ભાઈ}} — \emph{temno bhāī}, their brother
\end{itemize}

\begin{tcolorbox}[breakable,colback=teal!4,colframe=teal!35!black,title={Before you move on}]
What decides the last sign of \textbf{\gu{તેની} \gu{બહેન}}? (\textbf{\gu{બહેન}} is feminine.) Which \emph{our} takes the listener in? (\textbf{\gu{આપણું}}.)
\end{tcolorbox}
\section[Practice]{Yes or no — a practice}

\label{lesson:GU-C146-ask-practice}



\begin{quote}
Say \emph{her name is Mira}. Then \emph{our village} twice: once leaving the listener out, once taking the listener in.
\end{quote}

\subsection*{The exchange: likes, family and a plan}
\begin{quote}
\textbf{— \gu{શું} \gu{તમને} \gu{ચા} \gu{ગમે} \gu{છે}?} \textbf{— \gu{હા}, \gu{મને} \gu{ચા} \gu{ગમે} \gu{છે}. \gu{તમને}?} \textbf{— \gu{ના}, \gu{મને} \gu{ચા} \gu{નથી} \gu{ગમતી}. \gu{મને} \gu{દૂધ} \gu{ગમે} \gu{છે}.} \textbf{— \gu{તમારી} \gu{બહેન} \gu{અહીં} \gu{રહે} \gu{છે}, \gu{ખરું}?} \textbf{— \gu{હા}. \gu{તેનું} \gu{ઘર} \gu{અહીં} \gu{છે}.} \textbf{— \gu{તમે} \gu{આવો} \gu{છો} \gu{કે} \gu{નહીં}?} \textbf{— \gu{હા}, \gu{આવું} \gu{છું}!}
\end{quote}

Say these aloud:

\begin{itemize}
\raggedright
  \item the exchange aloud, taking each part in turn
  \item why it is \textbf{\gu{નથી} \gu{ગમતી}} and not \textbf{\gu{નથી} \gu{ગમતો}} (\textbf{\gu{ચા}} is feminine, and the tea is the subject)
\end{itemize}

\subsection*{Your turn}
\begin{itemize}
\raggedright
  \item \emph{Ask:} a friend whether she has a fever, with \textbf{\gu{શું}}
  \item \emph{Ask:} whether their village is far, with \textbf{\gu{ખરું}?} — \emph{temnũ gām dūr che, kharũ?}
  \item \emph{Say it:} \textbf{\gu{અહીં} \gu{કોઈ} \gu{નથી}, \gu{ખરું}?} — \emph{ahī̃ koī nathī, kharũ?}, nobody is here, right?
  \item \emph{Say it:} \textbf{\gu{શું} \gu{તમે} \gu{કંઈ} \gu{નથી} \gu{ખાતા}?} — \emph{śũ tame kaī̃ nathī khātā?}, aren't you eating anything?
\end{itemize}

\begin{tcolorbox}[breakable,colback=teal!4,colframe=teal!35!black,title={Before you move on}]
Which question asked for a choice? (\textbf{\gu{તમે} \gu{આવો} \gu{છો} \gu{કે} \gu{નહીં}?}) Which one expected \emph{yes}? (\textbf{…, \gu{ખરું}?})
\end{tcolorbox}
